\documentclass{article}
\usepackage{amsmath, amssymb, amsthm}
\newtheorem{theorem}{Theorem}[section]
\newtheorem{definition}[theorem]{Definition}
\newtheorem{example}[theorem]{Example}
\newtheorem{remark}[theorem]{Remark}

\begin{document}

\section*{Chapter 9: Euler's identities and Jacobi Triple Product Identity}

\subsection{Jacobi's Triple Product Identity}

Let $q$ and $z$ be complex numbers such that $|q| < 1$. One of the most elegant
identities involving $q$ and $z$ is perhaps Jacobi's Triple Product Identity given by
\begin{equation}
\sum_{j=-\infty}^{\infty} q^{j^2} z^j = \prod_{j=1}^{\infty} (1 + zq^{2j-1})(1 + z^{-1}q^{2j-1})(1 - q^{2j}).
\tag{9.1}
\end{equation}

There are several proofs of (9.1). In this chapter, we will present an elementary
proof due to G.E.\ Andrews. Before we proceed with the proof of (9.1), we give
a few interesting results that follow from this identity.

\begin{example}
Let $z = 1$ in (9.1). We immediately obtain a product representation
of the generating function for squares, namely,
\[
\sum_{j=-\infty}^{\infty} q^{j^2} = \prod_{j=1}^{\infty} (1 + q^{2j-1})^2(1 - q^{2j}).
\]
We remark here that the squares $j^2 = 1, 4$ and $9$ for $j = 1, 2$ and $3$ can be viewed
as the number of dots in a diagram.
\end{example}

\begin{example}
Let $z = q$ in (9.1), followed by replacing $q$ by $q^{1/2}$. We find that
\[
\sum_{j=-\infty}^{\infty} q^{j(j+1)/2} = 2\prod_{j=1}^{\infty} (1 + q^j)^2(1 - q^j).
\]
The integers $j(j+1)/2 = 1, 3, 6$ for $j = 1, 2$ and $3$ can be viewed as the number
of dots in a diagram.
\end{example}

\begin{example}
Let $q$ be replaced by $q^{3/2}$, followed by letting $z = -q^{1/2}$. We
immediately deduce from (9.1) that
\begin{equation}
\sum_{j=-\infty}^{\infty} (-1)^j q^{j(3j+1)/2} = \prod_{j=1}^{\infty} (1 - q^{3j-1})(1 - q^{3j-2})(1 - q^{3j}) = \prod_{j=1}^{\infty} (1 - q^j),
\tag{9.2}
\end{equation}
where the last equality follows from the fact that a positive integer must be of the
form $3j - r$, with $r = 0, 1$ or $2$. Note that by replacing $j$ by $-j$, we may write
the left hand of the above series as
\[
\sum_{j=-\infty}^{\infty} (-1)^j q^{j(3j-1)/2}.
\]
The integers of the form $\omega(j) = j(3j-1)/2$, with $j \geq 1$, are known as pentagonal
numbers. For $j = 1, 2$ and $3$, $\omega(j) = 1, 5$ and $12$.
\end{example}

\subsection{The terminating $q$-binomial Theorem}

Let $n \geq 1$ be an integer. The binomial theorem states that
\begin{equation}
(1 + x)^n = \sum_{k=0}^{n} \binom{n}{k} x^k
\tag{9.3}
\end{equation}
where
\begin{equation}
\binom{n}{k} = \frac{n!}{k!(n-k)!}.
\tag{9.4}
\end{equation}

There is a $q$-analogue of (9.3) and to describe this analogue, we first define the
$q$-analogue of (9.4).

\begin{definition}
Let $(q)_0 = 1$ and let
\[
(q)_n = \prod_{j=1}^{n} (1 - q^j)
\]
and
\begin{equation}
\binom{n}{k}_q = \frac{(q)_n}{(q)_k (q)_{n-k}}.
\tag{9.5}
\end{equation}
We also let
\[
(q)_\infty = \prod_{k=1}^{\infty} (1 - q^k).
\]
\end{definition}

The expression in (9.5) is a $q$-analogue of (9.4) since
\[
\lim_{q \to 1^-} \binom{n}{k}_q = \lim_{q \to 1^-} \frac{(1-q)(1-q^2)\cdots(1-q^n)}{(1-q)^n} \cdot \frac{(1-q)^k}{(q)_k} \cdot \frac{(1-q)^{n-k}}{(q)_{n-k}} = \binom{n}{k}.
\]

Observe further that the expressions for (9.4) and (9.5) motivate us to view the
$q$-analogue of $n!$ as $(q)_n$, even though $(q)_n$ does not approach $n!$ as $q$ approaches
$1^-$.

The $q$-binomial coefficient has several interesting properties. For example, one
can verify that
\begin{equation}
\binom{n}{k}_q = q^k \binom{n-1}{k}_q + \binom{n-1}{k-1}_q
\tag{9.6}
\end{equation}
and
\begin{equation}
\binom{n}{k}_q = \binom{n-1}{k}_q + q^{n-k} \binom{n-1}{k-1}_q.
\tag{9.7}
\end{equation}

When $q$ approaches $1^-$, both (9.6) and (9.7) reduce to the well known relation
\[
\binom{n}{k} = \binom{n-1}{k} + \binom{n-1}{k-1}.
\]

By (9.6), (9.7) and mathematical induction, we deduce that $\binom{n}{k}_q$ is a polynomial
in $q$. This is analogous to the fact that $\binom{n}{k}$ is an integer.

\begin{theorem}
Let $n$ be a positive integer, $x$ and $q$ be independent variables. Then
\begin{equation}
(1+x)(1+xq)\cdots(1+xq^{n-1}) = \sum_{k=0}^{n} \binom{n}{k}_q q^{k(k-1)/2} x^k.
\tag{9.8}
\end{equation}
\end{theorem}

\begin{proof}
We follow the approach of G.\ P\'olya and G.L.\ Alexanderson. Denote
$f(x) = (1+x)(1+xq)\cdots(1+xq^{n-1})$.
Note that
\[
(1+x)f(xq) = (1+x)(1+xq)\cdots(1+xq^{n-1})(1+xq^n) = f(x)(1+xq^n).
\]
If we write $f(x) = \sum_{k=0}^{n} Q_k x^k$,
where $Q_k$ is a function of $q$, then
\[
(1+x)\left(\sum_{k=0}^{n} q^k Q_k x^k\right) = (1+xq^n)\left(\sum_{k=0}^{n} Q_k x^k\right).
\]
Comparing the coefficient of $x^k$ on both sides, we deduce that
\[
q^k Q_k + q^{k-1} Q_{k-1} = Q_k + q^n Q_{k-1},
\]
which implies that
\[
(1-q^k) Q_k = q^{k-1}(1-q^{n-k+1}) Q_{k-1}.
\]
Hence, we find that
\begin{equation}
Q_k = q^{k-1} \frac{(1-q^{n-k+1})}{(1-q^k)} Q_{k-1}.
\tag{9.9}
\end{equation}
Observing that $Q_0 = 1$ and iterating (9.9), we deduce that
\[
Q_k = q^{(k-1)+(k-2)+\cdots+2+1} \frac{(1-q^{n-k+1})(1-q^{n-k+2})\cdots(1-q^n)}{(1-q^k)(1-q^{k-1})\cdots(1-q^2)(1-q)} Q_0 = q^{k(k-1)/2} \binom{n}{k}_q.
\]
This completes the proof of (9.8).
\end{proof}

\subsection{Two identities of Euler}

In this section, we derive two identities due to L.\ Euler. These identities may be
viewed as analogues of the identity
\[
e^x = \sum_{k=0}^{\infty} \frac{x^k}{k!}.
\]

We first rewrite (9.8) as
\[
(1+x)(1+xq)\cdots(1+xq^{n-1}) = \sum_{k=0}^{n} \frac{(q)_n}{(q)_k (q)_{n-k}} q^{k(k-1)/2} x^k.
\]
We may let $n \to \infty$ and deduce that
\begin{equation}
\prod_{k=1}^{\infty} (1 + xq^{k-1}) = \sum_{k=0}^{\infty} \frac{q^{k(k-1)/2}}{(q)_k} x^k.
\tag{9.10}
\end{equation}
This identity is due to Euler. Note that if we take $(q)_n$ as the analogue of $n!$ and
allow the numerator $q^{k(k-1)/2}$ to approach $1$, we could view the right hand side
of (9.10) as an analogue of the power series expansion of $e^x$.

Next, set $x = -1$ in (9.8) and multiply both sides by $(-1)^n/(q)_n$ to deduce that
\begin{equation}
0 = \sum_{k=0}^{n} \frac{(-1)^{n-1}}{(q)_{n-k}} \cdot \frac{q^{k(k-1)/2}}{(q)_k}, \quad n \geq 1.
\tag{9.11}
\end{equation}

Note that the right hand side is of the form $\sum_{k=0}^{n} a_k b_{n-k}$, which is the coefficient
of the $x^k$ in the expansion of $A(x)B(x)$ where
\[
A(x) = \sum_{m=0}^{\infty} \frac{(-1)^m x^m}{(q)_m} \quad\text{and}\quad B(x) = \sum_{m=0}^{\infty} \frac{q^{m(m-1)/2}}{(q)_m} x^m.
\]
By (9.11), we have
\begin{equation}
A(x)B(x) = 1.
\tag{9.12}
\end{equation}
Using (9.10), we find that
\begin{equation}
\sum_{m=0}^{\infty} \frac{(-1)^m x^m}{(q)_m} = \prod_{k=1}^{\infty} \frac{1}{(1+xq^{k-1})}.
\tag{9.13}
\end{equation}
Identity (9.13) is also due to Euler. By taking $(q)_n$ to be the analogue of $n!$, we
find that the left hand side of (9.13) is the analogue of the power series expansion
of $e^{-x}$. This shows that (9.12) is the analogue of $e^x \cdot e^{-x} = 1$.

\subsection{Andrews' proof of Jacobi's Triple Product Identity}

In this section, we will present Andrews' proof of Jacobi's Triple Product Identity.
We first replace $q$ by $q^2$ in (9.10), we find that
\begin{equation}
\prod_{k=1}^{\infty} (1 + xq^{2k-2}) = \sum_{k=0}^{\infty} \frac{q^{k(k-1)}}{(q^2)_k} x^k.
\tag{9.14}
\end{equation}
Replacing $x$ by $xq$ in (9.14), we find that
\begin{equation}
\prod_{k=1}^{\infty} (1 + xq^{2k-1}) = \sum_{k=0}^{\infty} \frac{q^{k^2}}{(q^2)_k} x^k.
\tag{9.15}
\end{equation}

Observing that
\[
\frac{1}{(q^2)_k} = \prod_{m=1}^{\infty} \frac{1-q^{2m+2k}}{1-q^{2m}},
\]
we rewrite (9.15) as
\[
\prod_{k=1}^{\infty} (1+xq^{2k-1}) = \sum_{k=0}^{\infty} \prod_{m=1}^{\infty} \frac{(1-q^{2m+2k})}{(1-q^{2m})} q^{k^2} x^k
= \prod_{m=1}^{\infty} \frac{1}{(1-q^{2m})} \sum_{k=0}^{\infty} \prod_{m=1}^{\infty} (1-q^{2m+2k}) q^{k^2} x^k
\]
\[
= \prod_{m=1}^{\infty} \frac{1}{(1-q^{2m})} \sum_{k=-\infty}^{\infty} \prod_{m=1}^{\infty} (1-q^{2m+2k}) q^{k^2} x^k,
\]
where in the last equality, we used the fact that for $k < 0$,
$\prod_{m=1}^{\infty} (1-q^{2m+2k}) = 0$.
This implies that
\begin{equation}
\prod_{k=1}^{\infty} (1+xq^{2k-1})(1-q^{2k}) = \sum_{k=-\infty}^{\infty} \prod_{m=1}^{\infty} (1-q^{2m+2k}) q^{k^2} x^k.
\tag{9.16}
\end{equation}

Next, using (9.14) with $x = -q^{2k+2}$, we find that
\[
\prod_{m=1}^{\infty} (1-q^{2m+2k}) = \sum_{\ell=0}^{\infty} \frac{(-1)^\ell q^{\ell^2+\ell+2k\ell}}{(q^2)_\ell}.
\]

Substituting this expression into the right hand side of (9.16), we deduce that
\[
\sum_{k=-\infty}^{\infty} q^{k^2} x^k \sum_{\ell=0}^{\infty} \frac{(-1)^\ell q^{\ell^2+\ell+2k\ell}}{(q^2)_\ell}
= \sum_{\ell=0}^{\infty} \sum_{k=-\infty}^{\infty} \frac{(-q)^\ell q^{(k+\ell)^2}}{(q^2)_\ell} x^k
\]
\[
= \sum_{\ell=0}^{\infty} \left(\sum_{k=-\infty}^{\infty} q^{(k+\ell)^2} x^{k+\ell}\right) \frac{(-q/x)^\ell}{(q^2)_\ell}
= \left(\sum_{k=-\infty}^{\infty} q^{k^2} x^k\right)\left(\sum_{\ell=0}^{\infty} \frac{(-q/x)^\ell}{(q^2)_\ell}\right)
\]
\begin{equation}
= \left(\sum_{k=-\infty}^{\infty} q^{k^2} x^k\right) \prod_{m=1}^{\infty} \frac{1}{(1+x^{-1}q^{2m-1})},
\tag{9.17}
\end{equation}
where we have used
\begin{equation}
\sum_{\ell=0}^{\infty} \frac{(-q/x)^\ell}{(q^2)_\ell} = \prod_{m=1}^{\infty} \frac{1}{(1+x^{-1}q^{2m-1})}
\tag{9.18}
\end{equation}
in the last equality. Identity (9.18) follows from (9.13) with $q$ replaced by $q^2$,
followed by replacing $x$ by $-q/x$. Combining (9.16) and (9.17), we conclude the
proof of Jacobi's Triple Product Identity.

\begin{remark}
Andrews' proof was motivated by Bellman's remark that there
was no elementary proof of the Jacobi Triple Product identity. Andrews' proof
is perhaps one of the first elementary proofs of the identity. Around the same
time, P.K.\ Menon independently discovered a proof similar to that of Andrews.
\end{remark}

\subsection{Number of representations of $n$ as sums of two squares}

We already know that $p \equiv 1 \pmod{4}$ if and only if $p$ is a sum of two squares.
One way to see this explicitly is to use the formula
\begin{equation}
r_2(n) = 4\left(\sum_{\substack{d \mid n \\ d \equiv 1 \pmod{4}}} 1 - \sum_{\substack{d \mid n \\ d \equiv 3 \pmod{4}}} 1\right),
\tag{9.19}
\end{equation}
where $r_2(n)$ is the number of ways of writing $n$ as a sum of two squares. Now if
$p \equiv 1 \pmod{4}$ then $r_2(p) = 8$. If $p \equiv 3 \pmod{4}$ then $r_2(p) = 0$.

In this section, we present M.D.\ Hirschhorn's proof of (9.19). Replace $z$ by
$-x^2 q$ and $q^2$ by $q$ in (9.1) to deduce that
\begin{equation}
(x - x^{-1})\prod_{k=1}^{\infty}(1-x^2 q^k)(1-x^{-2}q^k)(1-q^k) = \sum_{k=-\infty}^{\infty}(-1)^k x^{2k+1} q^{k(k+1)/2}.
\tag{9.20}
\end{equation}

We now set $k = 2n$ and $k = 2n+1$ in the sum on the right hand side and deduce that
\begin{align*}
(x - x^{-1})&\prod_{k=1}^{\infty}(1-x^2 q^k)(1-x^{-2}q^k)(1-q^k) \\
&= \sum_{n=-\infty}^{\infty} x^{4n+1} q^{2n^2+n} - \sum_{n=-\infty}^{\infty} x^{4n-1} q^{2n^2-n} \\
&= x\prod_{n=1}^{\infty}(1+x^4 q^{4n-1})(1+x^{-4}q^{4n-3})(1-q^{4n}) \\
&\quad - x^{-1}\prod_{n=1}^{\infty}(1+x^4 q^{4n-3})(1+x^{-4}q^{4n-1})(1-q^{4n}),
\end{align*}
where the last equality is obtained using (9.1).

Differentiating the last equality with respect to $x$ and setting $x = 1$, we deduce that
\begin{equation}
\prod_{n=1}^{\infty}(1-q^n)^3 = \prod_{n=1}^{\infty}(1+q^{4n-3})(1+q^{4n-1})(1-q^{4n})
\tag{9.21}
\end{equation}
\[
\times \left(1 - 4\sum_{n=1}^{\infty}\left(\frac{q^{4n-3}}{1+q^{4n-3}} - \frac{q^{4n-1}}{1+q^{4n-1}}\right)\right).
\]

This implies that (exercise)
\begin{equation}
\prod_{n=1}^{\infty}\frac{(1-q^{2n})(1-q^{2n-1})^2}{2} = 1 - 4\sum_{n=1}^{\infty}\left(\frac{q^{4n-3}}{1+q^{4n-3}} - \frac{q^{4n-1}}{1+q^{4n-1}}\right).
\tag{9.22}
\end{equation}

Note that by replacing $q$ by $-q$, the left hand side of (9.22) can be identified
using (9.1) as
\[
\left(\sum_{k=-\infty}^{\infty} q^{k^2}\right)^2
\]
which is $\sum_{n=0}^{\infty} r_2(n) q^n$.

Comparing the coefficients of the last series with the right hand side of (9.22)
(after replacing $q$ by $-q$), we arrive at (9.19).

\subsection{The partition function $p(n)$}

In this section, we will illustrate the use of Jacobi Triple Product Identity in the
study of the partition function $p(n)$.

\begin{definition}
A partition of a positive integer $n$ is a representation of $n$ as
a sum of non-decreasing positive integers. The partition function $p(n)$ is defined
as the number of partitions of $n$.
\end{definition}

\begin{example}
The integer 4 has the following partitions:
\[
4 = 1 + 3 = 2 + 2 = 1 + 1 + 2 = 1 + 1 + 1 + 1.
\]
The number of partitions of 4 is 5 and we write $p(4) = 5$.
\end{example}

Let $F(q) = \sum_{n=0}^{\infty} p(n) q^n$, where $p(0) := 1$. One of the most important identities
associated with $p(n)$ is the following:

\begin{theorem}
For $|q| < 1$, we have
\begin{equation}
\prod_{k=1}^{\infty} \frac{1}{(1-q^k)} = \sum_{n=0}^{\infty} p(n) q^n.
\tag{9.23}
\end{equation}
\end{theorem}

\begin{proof}
First let $q$ be real and that $0 < q < 1$. Consider the product
$F_m(q) = \prod_{k=1}^{m} \frac{1}{(1-q^k)}$.
If we expand the individual term in the product, we find that
\[
\prod_{k=1}^{m} \frac{1}{(1-q^k)} = (1 + q + q^{1+1} + q^{1+1+1} + \cdots)(1 + q^2 + q^{2+2} + q^{2+2+2} + \cdots)\cdots
(1 + q^m + q^{m+m} + \cdots).
\]
We now expand the products on the right hand side and denote the power series as
$\sum_{n=0}^{\infty} p_m(n) q^n$,
where $p_m(0) = 1$. For a fixed $n \geq 1$, we note that $p_m(n)$ counts the number
of representations of $n$ as a partition with parts not exceeding $m$.

Now, observe that $p_m(n) \leq p(n)$ and if $n \leq m$, $p_m(n) = p(n)$.
When $m$ approaches $\infty$, $p_m(n)$ approaches $p(n)$. Therefore,
$1 + \sum_{n=1}^{m} p(n) q^n \leq F_m(q) \leq F(q)$,
where $F(q) = \prod_{k=1}^{\infty} \frac{1}{(1-q^k)}$.
Hence, $1 + \sum_{n=1}^{\infty} p(n) q^n$ converges for $0 < q < 1$ uniformly in $m$. Now,
$\sum_{n=0}^{\infty} p_m(n) q^n \leq \sum_{n=0}^{\infty} p(n) q^n \leq F(q)$.
Hence, $F_m(q) \leq \sum_{n=0}^{\infty} p(n) q^n \leq F(q)$
and by letting $m \to \infty$, we complete the proof of the theorem.

We then extend the above identity by analytic continuation to the unit disk $|q| < 1$.
\end{proof}

Now, observe from Theorem 9.5 that
$1 = \left(\sum_{n=0}^{\infty} p(n) q^n\right) \cdot \left(\prod_{k=1}^{\infty}(1-q^k)\right)$.
By Example 9.3, we conclude that
\[
\left(\sum_{n=0}^{\infty} p(n) q^n\right) \cdot \left(1 + \sum_{n=1}^{\infty}(-1)^n(q^{\omega(n)} + q^{\omega(-n)})\right) = 1
\]
where $\omega(n) = n(3n-1)/2$. Hence we deduce that for $N \geq 1$,
\[
\sum_{n+\omega(\pm m)=N} p(n)(-1)^m = 0.
\]
In other words,
\[
p(N) = \sum_{\substack{n+\omega(\pm m)=N \\ n \neq N}} (-1)^{m+1} p(n) = \sum_{m=1}^{N-1}(-1)^{m+1} p(N - \omega(\pm m)).
\]

Major MacMahon tabulated $p(n)$ up to $n = 200$ using the above recurrence
formula for $p(n)$. From this table, S.\ Ramanujan observed the following amazing
congruences
\begin{align}
p(5n+4) &\equiv 0 \pmod{5} \tag{9.24} \\
p(7n+5) &\equiv 0 \pmod{7} \tag{9.25} \\
p(11n+6) &\equiv 0 \pmod{11}. \tag{9.26}
\end{align}

In the next section, we will give a proof of (9.24).

We end this section by giving another formula for computing $p(n)$. This formula
relates $p(n)$ and $\sigma(n)$ defined by $\sigma(n) = \sum_{\ell \mid n} \ell$.

By logarithmic differentiating (9.23), we find that
\[
\sum_{n=0}^{\infty} n\, p(n) q^n = \left(\sum_{n=1}^{\infty} \frac{nq^n}{1-q^n}\right) \cdot \left(\sum_{n=0}^{\infty} p(n) q^n\right),
\]
where we have used the fact that
$\sum_{n=1}^{\infty} \frac{nq^n}{1-q^n} = \sum_{n=1}^{\infty}\sum_{m=1}^{\infty} nq^{mn} = \sum_{s=1}^{\infty}\sigma(s) q^s$.

This immediately gives the recurrence formula for $p(n)$, namely,
\[
n\, p(n) = \sum_{\ell + m = n} \sigma(\ell) p(m).
\]

Next, instead of using (9.23), we start with
$G(q) = \prod_{n=1}^{\infty}(1-q^n)$,
then
\begin{equation}
q\frac{G'(q)}{G(q)} = \sum_{n=1}^{\infty}\sigma(n) q^n.
\tag{9.27}
\end{equation}

Using (9.3), we may write
$G'(q) = 1 + \sum_{\ell=1}^{\infty}(-1)^\ell \omega(\pm\ell) q^{\omega(\pm\ell)}$.

Furthermore, by using (9.23), we conclude from (9.27) that
\begin{equation}
\sigma(n) = -\sum_{\ell=0}^{\infty}(-1)^\ell \frac{\ell(3\ell+1)}{2}\, p\!\left(n - \frac{\ell(3\ell+1)}{2}\right)
\tag{9.28}
\end{equation}
\[
- \sum_{\ell=0}^{\infty}(-1)^{\ell+1}\frac{(\ell+1)(3\ell+2)}{2}\, p\!\left(n - \frac{(\ell+1)(3\ell+2)}{2}\right).
\]

The above formula allows us to compute $\sigma(n)$ even if we do not know the prime
factorization of $n$. In particular, if $\sigma(n) = n+1$ then $n$ is a prime.

\begin{example}
When $n = 19$, we note that $p(17) = 297$, $p(12) = 77$, $p(4) = 5$,
$p(18) = 385$, $p(14) = 135$ and $p(7) = 15$. By (9.28), we conclude that
\[
\sigma(19) = 2 \cdot 297 - 7 \cdot 77 + 15 \cdot 5 + 385 - 5 \cdot 135 + 12 \cdot 15 = 20.
\]
Note that $\sigma(19) = 20$ and hence we conclude that 19 is a prime.
\end{example}

We have shown above that values of $p(n)$ are useful in testing if an integer $N$
were a prime. It is therefore very useful if one has an explicit formula for $p(n)$.
Such a formula which arises from a divergent series was first discovered by G.H.\
Hardy and S.\ Ramanujan around 1918 and independently by J.V.\ Uspensky in
1920. In 1937, H.\ Rademacher discovered an exact formula of $p(n)$ in terms of a
convergent series. A.\ Selberg had also independently discovered the same formula
around the same time.

\subsection{Proof of $p(5n+4) \equiv 0 \pmod{5}$}

In this section, we give Ramanujan's proof of (9.24). We will need another identity
of Jacobi.

First, recall (9.20), namely,
\[
(x - x^{-1})\prod_{k=1}^{\infty}(1-x^2 q^k)(1-x^{-2}q^k)(1-q^k) = \sum_{k=-\infty}^{\infty}(-1)^k x^{2k+1} q^{k(k+1)/2}.
\]
We write the right hand side of the above identity as
$\sum_{k=0}^{\infty}(-1)^k(x^{2k+1} - x^{-(2k+1)}) q^{k(k+1)/2}$.
Dividing both sides of the identity by $(x - x^{-1})$ and letting $x$ tends to 1, we find that
\begin{equation}
\prod_{n=1}^{\infty}(1-q^n)^3 = \sum_{n=0}^{\infty}(-1)^n(2n+1) q^{n(n+1)/2}.
\tag{9.29}
\end{equation}

Ramanujan's original proof for the congruence for $p(5n+4)$ depends on (9.2) and (9.29). Write
\[
q\prod_{n=1}^{\infty}(1-q^n)^4 = \sum_{r=-\infty}^{\infty}\sum_{s=0}^{\infty}(-1)^{r+s}(2s+1) q^{1+r(3r+1)/2+s(s+1)/2}.
\]
When $1 + r(3r+1)/2 + s(s+1)/2$ is a multiple of 5, we find that
$2(r+1)^2 + (2s+1)^2 \equiv 0 \pmod{5}$.
But this can only happen when $2s+1$ is divisible by 5. Therefore, the coefficient
$a_{5n}$ of the series expansion
$q\prod_{n=1}^{\infty}(1-q^n)^4 = \sum_{k=0}^{\infty} a_k q^k$
is divisible by 5.

Hence, the coefficient of $q^{5N}$ of
$q\prod_{k=1}^{\infty}\frac{(1-q^{5k})}{(1-q^k)}$
is also $0 \pmod{5}$ and therefore,
\[
p(5n+4) \equiv 0 \pmod{5}.
\]

\end{document}
